\documentclass[11pt]{ctexart}
\usepackage[a4paper,top=2cm,bottom=2cm,left=2.2cm,right=2.2cm]{geometry}
\usepackage{fancyvrb}
\usepackage{xcolor}
\usepackage{fontspec}
\usepackage[most]{tcolorbox}
\usepackage{booktabs}
\usepackage{enumitem}
\usepackage{titlesec}
\usepackage{hyperref}

\setmonofont{Consolas}
\setCJKmonofont{SimHei}
\setlength{\parindent}{0pt}
\setlist{nosep,leftmargin=1.4em}

\definecolor{deep}{RGB}{22,50,92}
\definecolor{accent}{RGB}{26,116,88}
\definecolor{cmt}{RGB}{124,134,150}
\definecolor{core}{RGB}{192,86,24}
\definecolor{boxbg}{RGB}{250,251,253}
\definecolor{warm}{RGB}{255,248,240}
\definecolor{keybg}{RGB}{240,248,255}
\definecolor{keyline}{RGB}{100,150,200}
\definecolor{warnbg}{RGB}{255,243,235}
\definecolor{warnline}{RGB}{200,120,60}
\definecolor{okbg}{RGB}{240,252,244}
\definecolor{okline}{RGB}{26,116,88}
\definecolor{tipbg}{RGB}{248,248,252}
\definecolor{tipline}{RGB}{140,140,180}

\titleformat{\section}{\large\bfseries\color{deep}}{\thesection}{0.6em}{}
\titleformat{\subsection}{\normalsize\bfseries\color{accent}}{\thesubsection}{0.5em}{}
\titlespacing{\section}{0pt}{14pt}{6pt}
\titlespacing{\subsection}{0pt}{10pt}{4pt}

% 常用盒子
\newtcolorbox{keybox}[1]{enhanced, colback=keybg, colframe=keyline,
  boxrule=0.6pt, arc=1.5mm, left=3mm, right=3mm, top=2mm, bottom=2mm,
  title={\bfseries #1}, fonttitle=\small\bfseries, coltitle=white,
  colbacktitle=keyline}
\newtcolorbox{warnbox}[1]{enhanced, colback=warnbg, colframe=warnline,
  boxrule=0.6pt, arc=1.5mm, left=3mm, right=3mm, top=2mm, bottom=2mm,
  title={\bfseries #1}, fonttitle=\small\bfseries, coltitle=white,
  colbacktitle=warnline}
\newtcolorbox{okbox}[1]{enhanced, colback=okbg, colframe=okline,
  boxrule=0.6pt, arc=1.5mm, left=3mm, right=3mm, top=2mm, bottom=2mm,
  title={\bfseries #1}, fonttitle=\small\bfseries, coltitle=white,
  colbacktitle=okline}
\newtcolorbox{tipbox}[1]{enhanced, colback=tipbg, colframe=tipline,
  boxrule=0.5pt, arc=1.5mm, left=3mm, right=3mm, top=1.5mm, bottom=1.5mm,
  title={\bfseries #1}, fonttitle=\small\bfseries, coltitle=white,
  colbacktitle=tipline}

\begin{document}

% ==================== 封面 ====================
\begin{titlepage}
\centering
\vspace*{2.5cm}
{\color{deep}\Huge\bfseries 视频链路基座}\\[6pt]
{\color{deep}\Large\bfseries 从零手搓 Step1 教程}\\[14pt]
{\large 2026 FPGA 创新设计赛道（安路）赛题二}\\[6pt]
{\large 视觉方向 · Step1\_video\_link}\\[20pt]

\begin{tcolorbox}[enhanced, colback=boxbg, colframe=accent, boxrule=0.8pt,
  arc=2mm, width=12cm, left=4mm, right=4mm, top=3mm, bottom=3mm]
\centering
{\small 读者：零基础队员\\
目标：理解并手搓视觉链路第一部分（Step1）\\
前置要求：会开电脑、装了 ModelSim 和 TD\\
预计学习时间：10～14 天（每天 2～3 小时）}
\end{tcolorbox}

\vfill
{\small\color{black!50} fpga-ptz-tracker 教程系列 · 2026-09-19}
\end{titlepage}

% ==================== 目录 ====================
\tableofcontents
\newpage

% ============================================================
\section{你在造什么}

\subsection{一句话}

\textbf{把摄像头拍到的像素"接住"，把测试图案"生成"出来，把不同时钟的信号"安全传递"。}

翻译成电路：计数器、比较器、小存储器、几根线。没有复杂数学，没有算法——这是整个项目里最简单的一步，也是最好的入门点。

\subsection{在整个系统中的位置}

摄像头拍到的画面要经过好几级处理才能到屏幕上。Step1 是这条路的"地基"——它不做人脸识别、不做跟踪，它只保证\textbf{像素数据能从 A 点流到 B 点而不丢、不错、不停}。

\begin{keybox}{Step1 的三个模块}
\begin{enumerate}
  \item \textbf{ser96to24}：把摄像头送来的"4 像素一包"拆成"1 像素一拍"（串化器）
  \item \textbf{tpg\_720p}：没有摄像头时，假装自己是摄像头，生成测试图案（测试源）
  \item \textbf{pulse\_sync + cdc\_4phase}：把信号安全地从一个时钟域送到另一个时钟域（跨域桥）
\end{enumerate}
\end{keybox}

\subsection{为什么要先做这个}

后面的检测链（Step2：颜色识别）、闭环控制（Step3：云台转动）、OSD 叠加（Step4：屏幕画框），全都建立在这条"像素流"之上。如果像素流本身就有问题（丢数据、顺序错、帧同步乱），后面全部白搭。

\begin{okbox}{学完 Step1 你能做到}
\begin{itemize}
  \item 解释像素流从摄像头到屏幕的完整路径
  \item 看懂 Verilog 里的 always/assign/reg/wire
  \item 在 ModelSim 里跑仿真并看波形
  \item 手写一个计数器、一个串化器、一个测试图案发生器
\end{itemize}
\end{okbox}

% ============================================================
\section{阶段 0：三个数字电路概念（1～2 天）}

你不需要学完整门数电课。只需要理解三个东西，就能看懂本项目 90\% 的代码。

\subsection{时钟（clock）}

一个不停 0→1→0→1 的方波信号。板上有个晶振每秒产生 5000 万个这样的方波（50MHz）。

\textbf{关键点：所有动作都在时钟"从 0 变 1 的那一瞬间"发生。}就像军队的口令——喊"齐步走"的瞬间所有人同时迈步。时钟就是那个口令。

\subsection{寄存器（D 触发器）}

一个"在时钟上升沿锁存输入值的 1-bit 存储单元"。

\begin{Verbatim}[frame=single,fontsize=\small]
时钟：  ──┐ ┌──┐ ┌──┐ ┌──┐ ┌──
          └─┘  └─┘  └─┘  └─┘
          ↑    ↑    ↑    ↑
        上升沿 上升沿 上升沿 上升沿

输入 d:   1    0    1    1    ...
输出 q:   旧值  1    0    1    ...
               ↑    ↑    ↑
     q 在每个上升沿"记住"当时的 d
\end{Verbatim}

Verilog 写法：
\begin{Verbatim}[frame=single,fontsize=\small]
always @(posedge clk) begin
    q <= d;    // 每个时钟上升沿，q 取 d 的值
end
\end{Verbatim}

\subsection{组合逻辑}

没有存储、没有时钟——输入变了输出\textbf{立刻}变（有微小延迟）。例如加法器、比较器、多路选择器。

\begin{Verbatim}[frame=single,fontsize=\small]
wire [7:0] sum;
assign sum = a + b;    // b 变了 sum 立刻跟着变（几纳秒后）
\end{Verbatim}

\begin{warnbox}{组合逻辑 vs 时序逻辑——一字之差，天壤之别}
\textbf{时序逻辑}（有时钟，用 \texttt{always @(posedge clk)}）：输出在时钟沿才变，有"记忆"。\\
\textbf{组合逻辑}（无时钟，用 \texttt{assign}）：输出立刻跟输入变，无记忆。\\[4pt]
混淆的后果：仿真可能对但上板行为完全不同——这是新手最常踩的坑。
\end{warnbox}

\subsection{练习：写你的第一个 Verilog 计数器}

\begin{Verbatim}[frame=single,fontsize=\small]
// my_counter.v —— 一个 4 位计数器，0~15 循环
module my_counter (
    input  wire       clk,      // 时钟
    input  wire       rst_n,    // 复位（低电平有效）
    output reg  [3:0] count     // 4 位计数输出
);
    always @(posedge clk or negedge rst_n) begin
        if (!rst_n)
            count <= 4'd0;           // 复位时清零
        else
            count <= count + 4'd1;   // 每拍加 1
    end
endmodule
\end{Verbatim}

\begin{tipbox}{怎么跑仿真}
\begin{enumerate}
  \item 创建一个文件 \texttt{my\_counter.v}，粘贴上面的代码
  \item 再创建一个测试文件 \texttt{tb.v}：
\begin{Verbatim}[fontsize=\small]
module tb;
    reg clk = 0, rst_n = 0;
    always #5 clk = ~clk;          // 每 5ns 翻转 = 100MHz
    my_counter u (.clk(clk), .rst_n(rst_n), .count());
    initial begin
        #20 rst_n = 1;             // 20ns 后释放复位
        #200 $finish;              // 跑 200ns
    end
endmodule
\end{Verbatim}
  \item 在命令行：
\begin{Verbatim}[fontsize=\small]
vlib work
vlog my_counter.v tb.v
vsim -c tb -do "run -all; quit -f"
\end{Verbatim}
  \item 或用 GUI：\texttt{vsim tb} → 点 Run → 加波形看 count 从 0 数到 15
\end{enumerate}
\end{tipbox}

% ============================================================
\section{阶段 1：Verilog 最小语法集（3～5 天）}

你不需要学完整 Verilog。以下是你看懂本项目所需的全部语法，大约 30 分钟过一遍。

\subsection{必须会的 8 个语法}

\begin{Verbatim}[frame=single,fontsize=\small,numbers=left]
// 1. 模块声明（一个电路块）
module 模块名 #(参数) (端口列表);
  ...
endmodule

// 2. 端口类型
input  wire [7:0]  data_in    // 8 位输入
output reg  [23:0] data_out   // 24 位输出（reg 表示在 always 里赋值）
output wire        ready      // 1 位输出（wire 表示在 assign 里赋值）
inout  wire        sda        // 双向（如 I2C）

// 3. 时序逻辑（有时钟）
always @(posedge clk or negedge rst_n) begin
    if (!rst_n)
        q <= 0;
    else
        q <= d;
end

// 4. 组合逻辑
assign sum = a + b;

// 5. always 组合逻辑（多个条件时用这个代替多条 assign）
always @(*) begin
    if (sel) result = a;
    else     result = b;
end

// 6. 位拼接（把多个信号拼在一起）
wire [15:0] combined = {high_byte, low_byte};   // {8bit, 8bit} = 16bit

// 7. 位选择
wire bit3 = combined[3];          // 取第 3 位
wire [3:0] nibble = combined[7:4]; // 取高 4 位

// 8. 常量写法
8'd42     // 8 位十进制 42
8'h2A     // 8 位十六进制 2A
4'b1010   // 4 位二进制 1010
\end{Verbatim}

\subsection{三条铁律（违反 = 仿真对但上板错）}

\begin{warnbox}{铁律 1：时序逻辑用 \texttt{<=}，组合逻辑用 \texttt{=}}
\begin{Verbatim}[fontsize=\small]
// 正确
always @(posedge clk) q <= d;         // 时序：<=
assign sum = a + b;                    // 组合：不需要符号

// 错误
always @(posedge clk) q = d;          // 时序用了 = → 竞争冒险
\end{Verbatim}
\end{warnbox}

\begin{warnbox}{铁律 2：每个寄存器都要有复位值}
\begin{Verbatim}[fontsize=\small]
// 正确：复位时给确定的值
always @(posedge clk or negedge rst_n) begin
    if (!rst_n) count <= 0;           // ← 复位值
    else        count <= count + 1;
end

// 错误：没有复位 → 上电后寄存器是 X（未知态）
// 本项目实际踩过：tpg 的 pres 计数器漏了复位 → 输出全是 X
\end{Verbatim}
\end{warnbox}

\begin{warnbox}{铁律 3：有符号运算必须显式声明}
\begin{Verbatim}[fontsize=\small]
wire signed [12:0] err = $signed(cx) - $signed(center);
// 缺 $signed → 当无符号 → 负数变成巨大正数 → 系统跑飞
// 本项目实际踩过：闭环模型跑飞，查了半天是无符号减法
\end{Verbatim}
\end{warnbox}

\subsection{练习：读懂本项目的一个模块}

打开 \texttt{03\_vision/step1/03\_project/rtl/tpg\_720p.v}，找到以下结构并在旁边标注它们对应上面 8 个语法中的哪一个：

\begin{itemize}
  \item 模块声明和参数（第 1～10 行）
  \item 内部寄存器声明（\texttt{beat\_x}, \texttt{y}, \texttt{pres}）
  \item 颜色生成函数（\texttt{px\_color}）
  \item 主时序逻辑块（\texttt{always @(posedge I\_clk ...)}）
  \item 输出赋值（\texttt{O\_tdata <= ...}）
\end{itemize}

% ============================================================
\section{阶段 2：视频流概念（2～3 天）}

\subsection{像素流长什么样}

想象一根水管，每秒流出 5500 万个像素。每个像素是 24 位（R 8bit + G 8bit + B 8bit）。像素排队通过，不能停、不能丢。

\subsection{三个"路标"信号——全项目的通用语言}

\begin{keybox}{valid / sof / eol——必须吃透}
\begin{itemize}
  \item \textbf{valid}：这拍的数据是真的（不是垃圾）。类比：传送带上有东西。
  \item \textbf{sof}（Start of Frame）：这拍是\textbf{一帧的第一个像素}（左上角）。类比：新的一页纸开始了。
  \item \textbf{eol}（End of Line）：这拍是\textbf{一行的最后一个像素}（最右边）。类比：这行写完了，换行。
\end{itemize}
\end{keybox}

\subsection{一帧的完整结构}

\begin{Verbatim}[frame=single,fontsize=\small]
tuser=1, px(0,0), px(1,0), ..., px(1279,0), tlast=1   ← 第 0 行（320 拍）
tuser=0, px(0,1), px(1,1), ..., px(1279,1), tlast=1   ← 第 1 行
...
tuser=0, px(0,719), ..., px(1279,719), tlast=1        ← 第 719 行

一帧 = 1280 × 720 = 921,600 个像素
帧率 60fps → 每秒 55,296,000 个像素流过
\end{Verbatim}

\begin{tipbox}{坐标系约定}
\begin{itemize}
  \item 原点 (0,0) 在\textbf{左上角}
  \item x 向右增长（0..1279）
  \item y 向下增长（0..719）
  \item 这与 LCD/HDMI 扫描方向一致
\end{itemize}
\end{tipbox}

\subsection{96bit 打包：4 像素一拍}

为了减少布线，像素不是 1 个 1 个传的，而是 4 个打包成一拍：

\begin{Verbatim}[frame=single,fontsize=\small]
I_tdata[95:72]  = 像素 0（最左边）  {R[7:0], G[7:0], B[7:0]}
I_tdata[71:48]  = 像素 1
I_tdata[47:24]  = 像素 2
I_tdata[23:0]   = 像素 3（最右边）

一帧 1280 像素 ÷ 4 = 320 拍/行
一帧 = 320 拍 × 720 行 = 230,400 拍
\end{Verbatim}

\subsection{练习：跑仿真看波形}

\begin{enumerate}
  \item 在 ModelSim 里跑 \texttt{tb\_tpg}
  \item 加波形：\texttt{add wave /tb\_tpg/u\_dut/*}
  \item 跑 1 帧的时间（约 13ms），找：
  \begin{itemize}
    \item \texttt{tuser} 什么时候变高？（应该在帧的最开头，只出现一次）
    \item \texttt{tlast} 什么时候变高？（应该每 320 拍出现一次）
    \item \texttt{tdata} 在变化吗？（应该在不同竖条之间切换）
  \end{itemize}
\end{enumerate}

% ============================================================
\section{阶段 3：精读 ser96to24——你的第一个模块（3～5 天）}

\subsection{它做什么（生活类比）}

\begin{itemize}
  \item \textbf{输入}：一辆卡车每 4 秒来一次，装 4 个箱子（96bit = 4 × 24bit 像素）
  \item \textbf{输出}：每秒送走 1 个箱子（24bit 像素）
  \item \textbf{本质}：一个"小仓库"（FIFO），卡车来了卸 4 个进去，出口每秒取 1 个
\end{itemize}

\subsection{为什么要这么做}

后面的处理链（HSV/阈值/平滑）每个时钟周期只处理 1 个像素。如果不拆，就要写 4 份处理逻辑——资源翻 4 倍，调试难度翻 10 倍。

\subsection{打开代码逐段读}

文件：\texttt{03\_vision/step1/03\_project/rtl/ser96to24.v}（约 80 行）

\subsection{第 1 段：端口声明}

\begin{Verbatim}[frame=single,fontsize=\small,numbers=left]
module ser96to24 (
    input  wire        I_clk,       // 时钟
    input  wire        I_rst_n,     // 复位
    input  wire [95:0] I_tdata,     // 输入：4 像素打包
    input  wire        I_tvalid,    // 输入有效
    input  wire        I_tuser,     // 帧起始标记
    input  wire        I_tlast,     // 行尾标记
    output wire        O_tready,    // 永远 =1（"我随时能收"）
    output reg         O_valid,     // 输出有效
    output reg         O_sof,       // 输出帧起始
    output reg         O_eol,       // 输出行尾
    output reg  [23:0] O_data       // 输出：1 像素
);
\end{Verbatim}

\begin{tipbox}{O\_tready 为什么恒 1？}
这是"窃流"语义——视觉分支是旁路，不影响主链路。上游爱发多少发多少，我们只管"听"。
\end{tipbox}

\subsection{第 2 段：小仓库（FIFO）}

\begin{Verbatim}[frame=single,fontsize=\small,numbers=left]
reg [23:0] mem [0:7];      // 8 个格子，每个 24bit
reg [2:0]  wp;              // 写指针（指向下一个空格子）
reg [2:0]  rp;              // 读指针（指向最老的格子）
reg [3:0]  cnt;             // 里面有几个
\end{Verbatim}

\begin{keybox}{FIFO 的本质}
FIFO = First In, First Out（先进先出）。就像排队买奶茶：
先来的人先买到。写指针往队尾加人，读指针从队头取人。
\end{keybox}

\subsection{第 3 段：写入逻辑}

\begin{Verbatim}[frame=single,fontsize=\small,numbers=left]
wire do_wr = I_tvalid;      // 输入有效就写入

always @(posedge I_clk) begin
    if (do_wr) begin
        mem[wp]     <= I_tdata[95:72];  // 像素 0
        mem[wp_n]   <= I_tdata[71:48];  // 像素 1
        mem[wp_n2]  <= I_tdata[47:24];  // 像素 2
        mem[wp_n3]  <= I_tdata[23:0];   // 像素 3
        wp <= wp + 4;                    // 写指针跳 4 格
        cnt <= cnt + 4;                  // 计数 +4
    end
end
\end{Verbatim}

\subsection{第 4 段：读出逻辑}

\begin{Verbatim}[frame=single,fontsize=\small,numbers=left]
wire do_rd = (cnt != 0);    // 里面有数据就读

always @(posedge I_clk) begin
    if (do_rd) begin
        O_valid <= 1;
        O_data  <= mem[rp];    // 取最老的
        O_sof   <= sof_q[rp];  // 如果这是帧首像素
        O_eol   <= eol_q[rp];  // 如果这是行尾像素
        rp <= rp + 1;
        cnt <= cnt - 1;
    end else begin
        O_valid <= 0;
    end
end
\end{Verbatim}

\subsection{第 5 段：为什么不会溢出}

\begin{keybox}{收支平衡}
\begin{itemize}
  \item 输入：每 4 拍来 1 个 beat（含 4 像素）→ 平均 1 像素/拍
  \item 输出：每拍最多读 1 像素 → 1 像素/拍
  \item 平衡！缓冲峰值 7 < 深度 8，不会溢出
  \item 如果输入全速（每拍 1 beat），缓冲会爆——所以源端必须按 1/4 占空发数据
\end{itemize}
\end{keybox}

\subsection{练习}

\begin{enumerate}
  \item 画出 ser96to24 的框图（输入 → FIFO → 输出）
  \item 在 ModelSim 跑 \texttt{tb\_ser96to24}，看波形确认：
  \begin{itemize}
    \item 输入 4 像素一拍 → 输出 1 像素一拍
    \item sof 出现在帧的第一个像素
    \item eol 出现在行的最后一个像素
  \end{itemize}
  \item 尝试：把 FIFO 深度从 8 改成 16（\texttt{mem [0:15]}，指针改 4bit），重新仿真
\end{enumerate}

% ============================================================
\section{阶段 4：读懂 tpg\_720p——测试图案发生器（2～3 天）}

\subsection{它做什么}

没有摄像头时，TPG 假装自己是摄像头。它内部就是几个计数器驱动一个"查色表"：

\begin{Verbatim}[frame=single,fontsize=\small]
x 计数器：0→1279 循环（列位置）
y 计数器：0→719 循环（行位置）
颜色 = f(x, y)  ← 一个函数，根据位置返回颜色
\end{Verbatim}

\subsection{它长什么样}

\begin{Verbatim}[frame=single,fontsize=\small]
┌────────────────────────────────────────────────┐
│ 灰1  │ 灰2  │ 灰3  │ 灰4  │ 蓝  │ 绿  │ 深绿│ 深蓝│
│      │      │      │      │     │     │     │     │
│  ┌──────────────┐                               │
│  │   红色方块   │  ← 这是目标跟踪的"目标"       │
│  └──────────────┘                               │
│      （方块位置由外部输入控制，可以移动）       │
└────────────────────────────────────────────────┘
     8 个竖条各 160px 宽，全部避开红色阈值
\end{Verbatim}

\subsection{为什么竖条要"避开红色"}

因为后面的检测链在找红色目标。如果背景条也有红色，检测就会误识别。所以背景色选了灰阶（无饱和度）、蓝色（色相 240）、绿色（色相 120）——都不是红色（色相 0）。

\subsection{精读要点}

\begin{Verbatim}[frame=single,fontsize=\small,numbers=left]
// 核心：一个函数根据 x 坐标返回颜色
function [23:0] px_color;
    if (在方块内)     px_color = 24'hE0_20_20;  // 红色
    else if (px < 160)  px_color = 24'h40_40_40;  // 灰条 1
    else if (px < 320)  px_color = 24'h80_80_80;  // 灰条 2
    ...
endfunction

// 每拍输出 4 个像素（一个 beat）
O_tdata <= {px_color(x×4+0, y), px_color(x×4+1, y),
            px_color(x×4+2, y), px_color(x×4+3, y)};
\end{Verbatim}

\subsection{练习}

\begin{enumerate}
  \item 改一种条的颜色（如把蓝色改成紫色）→ 重新仿真 → 确认输出变化
  \item 改方块位置（\texttt{I\_box\_x} 和 \texttt{I\_box\_y}）→ 确认方块移动
  \item 思考：如果把一个条改成红色（E02020），检测链会怎么反应？
\end{enumerate}

% ============================================================
\section{阶段 5：CDC 概念（可选，2～3 天）}

\subsection{什么是跨时钟域}

芯片里不同模块可能用不同的时钟。比如：摄像头数据用 74.25MHz，舵机控制用 50MHz。当一个信号从 74.25MHz 的模块传到 50MHz 的模块时，就会出问题——因为接收方可能在信号变化的"一半"时采样到不确定的值。

\begin{warnbox}{为什么不能直接连线}
如果 74.25MHz 信号在 50MHz 时钟上升沿的\textbf{同一时刻}变化，接收寄存器可能进入
"亚稳态"（既不是 0 也不是 1），后续电路全部出错。
\end{warnbox}

\subsection{两根"安全桥"}

\begin{keybox}{pulse\_sync（单 bit 脉冲跨域，12 行）}
原理：源域每发一个脉冲，翻转一个寄存器。目标域用两个寄存器同步这个翻转，再检测"变化"产生脉冲。
\end{keybox}

\begin{keybox}{cdc\_4phase（多 bit 数据跨域，80 行）}
原理：源域把数据放在总线上 + 翻转一个 req 信号。目标域发现 req 变化后，锁存数据，翻转 ack 回应。源域收到 ack 后可以发下一个。\\
就像寄快递：你把包裹放门口（数据+req），对方取走后放个签收回执（ack），你看到回执才放下一个包裹。
\end{keybox}

\subsection{学习建议}

先跳过这两个文件（不影响理解其他模块），等学完 Step2 的检测链后再回来看。只需要记住一条规则：\textbf{跨时钟域不能裸连，必须走这两个模块}。

% ============================================================
\section{完整验证：跑通 Step1 全部测试}

\subsection{一键回归}

\begin{Verbatim}[frame=single,fontsize=\small]
cd D:\fpga-ptz-tracker\tools
python run_all.py
\end{Verbatim}

期望输出：\texttt{汇总: 18/18 PASS}（仓库全量口径；本册阶段为 15 项，其后
Step5/闭环补强又新增 3 项）

其中 Step1 相关的 4 个测试：

\begin{itemize}
  \item \texttt{tb\_ser96to24}：随机数据 + 3 帧，队列自检（数据/sof/eol 逐像素）
  \item \texttt{tb\_tpg}：几何计数断言 + 整帧 921600 像素与金标准比对
  \item \texttt{tb\_pulse\_sync}：异步双钟 100 脉冲 1:1
  \item \texttt{tb\_cdc4phase}：异步双钟 200 随机数据完整性
\end{itemize}

\subsection{单跑某个仿真}

\begin{Verbatim}[frame=single,fontsize=\small]
cd 03_vision/step1_video_link/03_project/sim

# 编译
vlib work
vlog ../rtl/ser96to24.v ../tb/tb_ser96to24.v

# 跑仿真（命令行模式）
vsim -c tb_ser96to24 -do "run -all; quit -f"

# 跑仿真（GUI 模式，可以看波形）
vsim tb_ser96to24
# 在 GUI 里: add wave /*  →  Run All
\end{Verbatim}

% ============================================================
\part{工程设计详解：Step1 的实现分析}

\section{工程总览}

\subsection{文件清单}

Step1 的全部工程文件位于 \texttt{03\_vision/step1\_video\_link/03\_project/}：

\begin{center}
\renewcommand{\arraystretch}{1.25}
\begin{tabular}{llll}
\toprule
\textbf{文件} & \textbf{行数} & \textbf{类型} & \textbf{作用} \\
\midrule
ser96to24.v & 96 & 可综合 & 96bit 转 1 像素/拍串化器 \\
tpg\_720p.v & 105 & 可综合 & 720p 测试图案发生器 \\
pulse\_sync.v & 37 & 可综合 & 单 bit 脉冲跨时钟域 \\
cdc\_4phase.v & 78 & 可综合 & 多 bit 数据 4 相握手跨域 \\
sim\_video\_src.v & 73 & 仅仿真 & 回放文件向量 \\
frame\_dump24.v & 68 & 仅仿真 & 帧转储为 hex（供 Python 比对）\\
tb\_ser96to24.v & 90 & 测试 & 串化器 TB \\
tb\_tpg.v & 72 & 测试 & TPG 几何断言 + 帧比对 \\
tb\_pulse\_sync.v & 45 & 测试 & 跨域脉冲 1:1 \\
tb\_cdc4phase.v & 63 & 测试 & 跨域数据完整性 \\
\bottomrule
\end{tabular}
\end{center}

\subsection{数据流}

\begin{Verbatim}[frame=single,fontsize=\small]
                    ┌──────────────────────────────┐
 96bit AXI-S  ─────>│  ser96to24                   │─────> 1 像素/拍 24bit
 (4 像素/拍)        │  (拆包 + 8 深 FIFO)          │       (给检测链)
                    └──────────────────────────────┘

 无摄像头时，用 tpg_720p 替代上游：
                    ┌──────────────────────────────┐
 tpg_720p     ─────>│  ser96to24（同一模块）        │─────> 1 像素/拍
 (测试图案)         └──────────────────────────────┘

 跨时钟域（两个方向共用）:
   pulse_sync : 单 bit 脉冲（帧拍）        74MHz -> 50MHz
   cdc_4phase : 多 bit 数据（目标包）      74MHz -> 50MHz
\end{Verbatim}

\subsection{三条设计原则}

\begin{okbox}{贯穿全工程的设计原则}
\begin{enumerate}
  \item \textbf{窃流旁路}：视觉分支绝不反压主线（\texttt{O\_tready} 恒 1），
        出错可随时摘除，显示链路行为与官方例程完全一致
  \item \textbf{结构自对齐优先}：能用"读写同指针"解决对齐，就不用"移位+手工推相位"
        （后者在 Step2 的平滑模块上让团队返工了四版）
  \item \textbf{金标准驱动}：所有算术先在 Python 建位精确模型，RTL 逐位复刻并回归验证
\end{enumerate}
\end{okbox}

% ============================================================
\section{模块详解：ser96to24}

\subsection{为什么需要它}

摄像头/ISP 输出的是 4 像素打包（96bit/拍），这样布线少、带宽高。
但后续检测链（HSV 变换、阈值、平滑）每拍只能处理 1 个像素——
如果按 4 像素写，就要 4 份并行逻辑：资源 $\times$4、调试难度 $\times$10。

\textbf{设计决策：在链路入口做一次"降速"，把 4 像素/拍拆成 1 像素/拍。}

\subsection{核心结构：FIFO}

\begin{Verbatim}[frame=single,fontsize=\small,numbers=left]
// 8 深缓冲：数据与标志同址存放
reg [23:0] mem   [0:7];     // 像素数据
reg        sof_q [0:7];     // 帧首标志
reg        eol_q [0:7];     // 行尾标志
reg [2:0]  wp, rp;          // 写指针 / 读指针
reg [3:0]  cnt;             // 当前存了几个
\end{Verbatim}

\textbf{为什么深度是 8？}\quad
输入每 4 拍来 4 像素（$+4$），输出每拍取 1 像素（$-1$）。
最坏情况净积压 $+3$/拍；考虑行尾边界与读空闲，实测峰值深度 7。
深度 8 留 1 格余量，且 8 是 2 的幂——指针用 3bit 自然回绕，无需取模。

\subsection{写入逻辑（4 像素一次性入队）}

\begin{Verbatim}[frame=single,fontsize=\small,numbers=left]
wire do_wr = I_tvalid & O_tready;   // O_tready 恒 1，故等价于 I_tvalid

always @(posedge I_clk or negedge I_rst_n) begin
    if (!I_rst_n) begin
        wp <= 3'd0; rp <= 3'd0; cnt <= 4'd0;
    end else begin
        if (do_wr) begin
            mem[wp]      <= p0;         // 像素 0（最左）
            mem[wp_n]    <= p1;
            mem[wp_n2]   <= p2;
            mem[wp_n3]   <= p3;         // 像素 3（最右）
            sof_q[wp]    <= I_tuser;    // 帧首标志只挂第 1 个像素
            eol_q[wp_n3] <= I_tlast;    // 行尾标志只挂第 4 个像素
            wp <= wp + 4'd4;
        end
        ...
\end{Verbatim}

\begin{keybox}{关键点 1：标志与像素严格绑定}
\texttt{tuser} 表示"这一拍是帧首"，但展开后真正的帧首只有 4 个像素中的\textbf{第 1 个}。
所以 \texttt{sof\_q[wp]} 只给第 1 格，其余 3 格给 0。
\texttt{tlast} 同理，只给第 4 个像素。\\[4pt]
这一步如果写错（比如给整拍 4 个像素都置位），下游的帧计数会变成 4 倍。
\end{keybox}

\textbf{关键点 2：指针步进 4}\quad
一次写入占 4 格，所以 \texttt{wp += 4}。指针位宽 3bit（0\textasciitilde 7），
加 4 后自然回绕，无需额外逻辑。

\subsection{读出逻辑（每拍 1 像素）}

\begin{Verbatim}[frame=single,fontsize=\small,numbers=left]
wire do_rd = (cnt != 4'd0);     // 有数据就读

always @(*) begin
    O_valid = do_rd;            // 组合输出
    O_data  = mem[rp];
    O_sof   = sof_q[rp];
    O_eol   = eol_q[rp];
end
\end{Verbatim}

\textbf{为什么用组合读出而不是寄存输出？}\quad
分布式 RAM 的组合读延迟很短（几个 LUT），在 74MHz 下时序满足。
若改成寄存器输出，整个检测链会多延迟 1 拍，给帧内对齐带来额外复杂度。

\subsection{计数器更新（同时读写）}

\begin{Verbatim}[frame=single,fontsize=\small,numbers=left]
case ({do_wr, do_rd})
    2'b10:   cnt <= cnt + 4'd4;    // 只写
    2'b01:   cnt <= cnt - 4'd1;    // 只读
    2'b11:   cnt <= cnt + 4'd3;    // 同时读写：+4-1 = +3
    default: cnt <= cnt;
endcase
\end{Verbatim}

\begin{warnbox}{新手最容易写错的地方}
分开写两个 \texttt{if} 会出错：
\begin{Verbatim}[fontsize=\small]
if (do_wr) cnt <= cnt + 4;
if (do_rd) cnt <= cnt - 1;   // ← 错误！后一句覆盖前一句
\end{Verbatim}
同时发生时，后者把前者的结果覆盖掉，计数逐渐跑偏（每拍少加 4 或多减 1）。
\textbf{必须用 case 把四种组合合并处理}——这是 FIFO 的标准写法。
\end{warnbox}

\subsection{溢出保护（仅仿真）}

\begin{Verbatim}[frame=single,fontsize=\small]
// synthesis translate_off
always @(posedge I_clk) begin
    if (I_rst_n && (cnt > 4'd7))
        $display("SER96TO24_FATAL: overflow cnt=%0d t=%0t", cnt, $time);
end
// synthesis translate_on
\end{Verbatim}

\texttt{translate\_off/on} 之间的代码只在仿真时存在，综合时被完全忽略——
既能第一时间暴露问题，又不占任何硬件资源。

\begin{warnbox}{这个保护立刻抓到了真 bug}
首版仿真时 TB 和 TPG 都以\textbf{满速}发 beat（相当于 4 倍像素率），
断言立刻报 \texttt{overflow cnt=13}。\\[4pt]
\textbf{根因}：真实 AXI-S 总线上 beat 占空比是 1/4（beat 率 = 像素率/4），
不是每拍都发。720p60 的像素率 55.3Mpx/s ÷ 4 = 13.8M beat/s，
在 74.25MHz 下就是"每 4 拍 1 个 beat"。\\[4pt]
\textbf{修正}：TPG / TB / sim\_video\_src 统一加 4 分频节拍。
这个 bug 如果留到上板，现象会是"画面偶尔丢像素"，极难定位。
\end{warnbox}

% ============================================================
\section{模块详解：tpg\_720p}

\subsection{三个用途}

\begin{enumerate}
  \item \textbf{无摄像头演示}：板子上没插摄像头时也能出图，独立演示视觉链路
  \item \textbf{闭环仿真驱动}：方块位置由外部输入控制，TB 可用云台模型实时改变它（相机在环）
  \item \textbf{确定性测试源}：图案完全可预测，Python 能算出期望帧做逐像素比对
\end{enumerate}

\subsection{图案设计：为什么是这 8 个颜色}

\begin{center}
\renewcommand{\arraystretch}{1.2}
\begin{tabular}{llll}
\toprule
\textbf{竖条} & \textbf{RGB} & \textbf{HSV（近似）} & \textbf{为什么选它} \\
\midrule
灰 1 & 404040 & H=0, S=0, V=64 & 无饱和度，不触发任何颜色阈值 \\
灰 2 & 808080 & H=0, S=0, V=128 & 同上，验证中灰 \\
灰 3 & C0C0C0 & H=0, S=0, V=192 & 同上，验证亮灰 \\
灰 4 & E0E0E0 & H=0, S=0, V=224 & 高亮但无色——测"高亮无饱和"边界 \\
蓝   & 2020E0 & H=240 & 与红色（H=0）相差最远的色相 \\
绿   & 20E020 & H=120 & 三基色之一，验证中间色相 \\
深绿 & 106010 & H=120, V=96 & 低亮度绿色——测暗色 \\
深蓝 & 101060 & H=240, V=96 & 低亮度蓝色——测暗色 \\
\midrule
\textbf{红方块} & E02020 & H=0, S=192, V=224 & \textbf{唯一目标}，高饱和高亮度红 \\
\bottomrule
\end{tabular}
\end{center}

\textbf{设计要点：所有背景色都避开了红色阈值窗（H$\in$[335,25] 且 S$\geq$80 且 V$\geq$60）。}\quad
灰阶 S=0 不满足饱和度条件；蓝色 H=240、绿色 H=120 都不在色相窗内。
这样检测链一眼锁定红方块，不受背景干扰——这是后面"逐像素金标准比对"能通过的前提。

\subsection{几何计数器与颜色函数}

\begin{Verbatim}[frame=single,fontsize=\small,numbers=left]
localparam BEATS_PER_LINE = IMG_WIDTH / 4;   // 1280/4 = 320

reg [10:0] beat_x;   // 0..319（打包坐标）
reg [9:0]  y;        // 0..719（行坐标）
reg [1:0]  pres;     // 4 分频节拍

function [23:0] px_color;
    input [12:0] px;    // 真实像素列号
    input [9:0]  py;
    begin
        if (BOX_EN && (px >= I_box_x) && (px < I_box_x + BOX_W) &&
            (py >= I_box_y) && (py < I_box_y + BOX_H))
            px_color = 24'hE0_20_20;              // 红方块
        else if (px < 13'd160)  px_color = 24'h40_40_40;
        else if (px < 13'd320)  px_color = 24'h80_80_80;
        ...  // 共 8 段
    end
endfunction
\end{Verbatim}

\begin{warnbox}{这里曾经写错过一次}
首版用位切片判竖条：\texttt{case (px[12:8])}。
但 160 不是 2 的幂，位切片得到的分段与 160px 等宽完全不符。\\[4pt]
\textbf{修正}：改用比较链（\texttt{px < 160 / px < 320 / ...}）。
综合器会把比较链优化成并联比较器，资源几乎不变。
\end{warnbox}

\subsection{4 分频节拍（关键设计）}

\begin{Verbatim}[frame=single,fontsize=\small]
reg [1:0] pres;
wire do_send = (pres == 2'd0) & (~GAP_EN | ...);

always @(posedge I_clk) begin
    pres <= pres + 2'd1;          // 0,1,2,3,0,1,2,3...
    if (do_send) begin
        // 输出一个 beat（4 像素）
    end
end
\end{Verbatim}

\textbf{为什么必须 4 分频？}\quad
720p60 的像素率 55.3Mpx/s，在 74.25MHz 时钟下 beat 率只有 13.8M/s，
即\textbf{每 4 拍 1 个 beat}。如果 TPG 满速发，下游 FIFO 立刻溢出
（这就是上一节那个真 bug）。

\begin{warnbox}{踩过的坑：复位遗漏}
首版 \texttt{pres} 忘了写进复位分支，上电后是 X（未知态），
\texttt{pres == 2'd0} 永远不成立 $\to$ TPG 一个像素都不输出。
\\[4pt]
\textbf{教训：每个寄存器都必须有复位值。}这是新手最常犯的错误，
而且现象往往是"完全没输出"，很难联想到复位。
\end{warnbox}

\subsection{帧同步标志}

\begin{Verbatim}[frame=single,fontsize=\small]
O_tuser <= (beat_x == 0) && (y == 0);           // 整帧仅一拍为 1
O_tlast <= (beat_x == BEATS_PER_LINE - 1);      // 每行最后一拍
\end{Verbatim}

\texttt{tuser} 只在"第 0 行第 0 拍"为 1——即一帧的最开头；
\texttt{tlast} 在每行第 319 拍为 1——即一行的最右端。

\subsection{LFSR 随机空闲（可选，锻炼鲁棒性）}

\texttt{GAP\_EN=1} 时用线性反馈移位寄存器产生约 7\% 概率的空闲拍，
模拟真实总线的随机间隙：

\begin{Verbatim}[frame=single,fontsize=\small]
lfsr <= {lfsr[14:0], lfsr[15] ^ lfsr[13] ^ lfsr[12] ^ lfsr[10]};
wire do_send = (pres == 2'd0) && (~GAP_EN | (lfsr[0] | (lfsr[7:0] == 8'h00)));
\end{Verbatim}

% ============================================================
\section{模块详解：pulse\_sync（单 bit 脉冲跨域）}

\subsection{问题}

一个 74.25MHz 域的"单拍脉冲"（例如每帧一次的帧拍），要传给 50MHz 域。
直接连一根线会出两个问题：

\begin{itemize}
  \item 50MHz 采样时可能采到脉冲的\textbf{边沿} $\to$ 亚稳态
  \item 74.25MHz 的单拍脉冲只有 13.5ns 宽，50MHz 周期是 20ns——
        可能被\textbf{完全错过}（一次都采不到）
\end{itemize}

\subsection{解决方案：toggle + 边沿还原}

\begin{Verbatim}[frame=single,fontsize=\small,numbers=left]
// 源域：每来一个脉冲，翻转一次电平
reg toggle_src;
always @(posedge I_clk_src or negedge I_rst_src_n) begin
    if (!I_rst_src_n)     toggle_src <= 1'b0;
    else if (I_pulse_src) toggle_src <= ~toggle_src;
end

// 目标域：3 级同步 + 边沿检测
reg [2:0] sync_dst;
always @(posedge I_clk_dst or negedge I_rst_dst_n) begin
    if (!I_rst_dst_n) sync_dst <= 3'b000;
    else              sync_dst <= {sync_dst[1:0], toggle_src};
end

assign O_pulse_dst = sync_dst[1] ^ sync_dst[2];   // 电平变化 = 一个脉冲
\end{Verbatim}

\textbf{原理}\quad
电平（持续高或持续低）比单拍脉冲"宽"，容易安全采样。
目标域检测到电平发生翻转（\texttt{sync\_dst[1] != sync\_dst[2]}），
就知道"源域发了一个脉冲"，于是产生自己域内的一个单拍脉冲。

\begin{keybox}{约束条件（必须满足）}
源脉冲之间的间隔必须大于目标域 3 个时钟周期，
否则两次翻转来不及被同步，会漏脉冲。\\[4pt]
本项目用于"每帧一次"的帧拍——帧间隔 16.7ms，远大于 60ns，绝对安全。
\end{keybox}

% ============================================================
\section{模块详解：cdc\_4phase（多 bit 数据跨域）}

\subsection{问题}

目标包（68bit，拓展后 135bit）要从检测域传到显示域/云台域。
多 bit 不能像单 bit 那样简单同步——因为每根线的延迟略有差异，
接收方可能采到"新旧混合"的值（例如高 8 位是新坐标、低 8 位是旧坐标）。

\subsection{解决方案：4 相握手}

\begin{Verbatim}[frame=single,fontsize=\small]
源域                                      目标域
────                                       ────
1. 数据放总线上（保持不动）
2. 翻转 req  ──────────────────────────> 3. 检测到 req 变化
                                          4. 锁存数据（此时数据已稳定）
                                          5. 翻转 ack
6. 检测到 ack 变化 <──────────────────────┘
7. 解除 busy，可以发下一包
\end{Verbatim}

\textbf{为什么能保证数据正确？}\quad
关键是\textbf{数据在 req 翻转之前就已稳定，且整个握手期间保持不变}。
目标域在检测到 req 变化后的任意时刻采样，都采到完整的新数据。

\begin{keybox}{类比：寄快递签收}
你把包裹放门口（数据），贴一张"已放件"标签（req 翻转）。
快递员取走包裹，贴一张"已取件"回执（ack 翻转）。
你看到回执才放下一个包裹——绝不会出现"包裹只放了一半就被取走"。
\end{keybox}

\subsection{关键实现与一个真实 bug}

\begin{Verbatim}[frame=single,fontsize=\small,numbers=left]
// 源域：busy 的解除条件
if (!busy) begin
    if (I_src_load) begin
        data_hold <= I_src_data;   // 1. 先锁存数据
        req       <= ~req;         // 2. 再翻转 req
        busy      <= 1'b1;
    end
end else if (ack_2ff == req) begin
    busy <= 1'b0;                  // 3. 用"电平比较"判断对方已应答
end
\end{Verbatim}

\begin{warnbox}{真实 bug：TB 驱动的同沿竞态}
首版仿真发现\textbf{每隔一包丢一包}（1,3,5,7\ldots 收到，2,4,6 丢失）。\\[4pt]
\textbf{根因不在 RTL，在测试平台}：\texttt{wait(!busy)} 在 busy 被置位的\textbf{同一个时钟沿}
就通过了（因为非阻塞赋值此时还未生效），于是第二个包在 DUT 仍忙时发出，被丢弃。\\[4pt]
\textbf{修正}：驱动侧改为"先等 busy 上升 $\to$ 再等 busy 下降"的稳健时序。\\[4pt]
\textbf{教训}：跨域模块的 TB 必须显式等待握手完成，不能假设一个时钟沿就够。
定位这个 bug 用了逐拍波形追踪（打印 req/busy/ack 的每拍值），
最后靠"数据序列号"看出是隔包丢失——如果是随机数据根本看不出规律。
\end{warnbox}

\subsection{验证}

\texttt{tb\_cdc4phase} 用两个异步时钟（周期 14ns 与 20ns，无固定相位关系），
发送 200 个随机 23bit 数据，逐个校验完整性。全部通过。

% ============================================================
\section{仿真专用模块}

\subsection{sim\_video\_src（文件回放源）}

从 \texttt{chain\_beats.hex}（96bit 拍数据）与 \texttt{chain\_flags.hex}（标志位）读入
Python 生成的确定性场景，按帧重复播放。

\begin{warnbox}{踩过的坑：readmemh 的位宽限制}
首版 flags 文件每行写两个字符（如 \texttt{"10"} 表示 tuser=1, tlast=0），
但 TB 里声明的是 \texttt{reg [1:0] flags}——只允许 \textbf{1 位}十六进制数字，
readmemh 直接截断，所有标志位全毁（tuser/tlast 永远为 0）。\\[4pt]
\textbf{修正}：改成单字符格式（值 0\textasciitilde 3：\texttt{2}=tuser、\texttt{1}=tlast、\texttt{3}=两者）。\\[4pt]
\textbf{教训}：hex 向量文件的字段宽度 = 向量位宽/4，宁可浪费一位也不要溢出。
\end{warnbox}

\subsection{frame\_dump24（帧转储）}

把 vpix 流（1 像素/拍）写成 hex 文件，每行一个像素的 RGB888：

\begin{Verbatim}[frame=single,fontsize=\small]
$fwrite(fd, "%02X%02X%02X\n", I_data[23:16], I_data[15:8], I_data[7:0]);
\end{Verbatim}

\begin{tipbox}{小陷阱：ModelSim 的 \%X 输出小写}
Python 比对时要注意：\texttt{c0c0c0} 与 \texttt{C0C0C0} 是同一个颜色。
\texttt{run\_all.py} 里统一做了 \texttt{.upper()} 归一化，避免这类假失配。
\end{tipbox}

% ============================================================
\section{测试平台设计}

\subsection{四层验证策略}

\begin{center}
\renewcommand{\arraystretch}{1.25}
\begin{tabular}{lll}
\toprule
\textbf{层次} & \textbf{方法} & \textbf{能抓到的错误} \\
\midrule
模块级 & 随机向量 + 逐位比对 & 算术错误、标志错位 \\
帧级   & 整帧 921,600 像素比对 & 几何错误、边界处理错误 \\
协议级 & 异步双钟 + 随机数据 & 跨域丢包、亚稳态 \\
系统级 & 端到端场景 & 模块间接口不匹配 \\
\bottomrule
\end{tabular}
\end{center}

\subsection{tb\_ser96to24 的关键设计：唯一数据编码}

\begin{Verbatim}[frame=single,fontsize=\small,numbers=left]
// 期望队列：把输入展开后逐像素压入
for (k = 0; k < 4; k = k + 1) begin
    exp_d[push_n + k] = 24'h100000 * f + b * 4 + k;  // 唯一编码
    exp_s[push_n + k] = (b == 0) && (k == 0);         // 帧首只在第 1 像素
    exp_e[push_n + k] = (k == 3) && ((b % BEATS) == BEATS-1); // 行尾只在第 4 像素
end
\end{Verbatim}

\textbf{设计要点：数据用唯一编码}\quad
第 $f$ 帧、第 $b$ 拍、第 $k$ 个像素的值 = \texttt{0x100000*f + b*4 + k}。
每个像素值全网唯一——出错时一眼看出是\textbf{哪个像素}错位了，
而不只是知道"有错误"。这比随机数据强大得多。

\subsection{tb\_tpg 的双重验证}

\begin{enumerate}
  \item \textbf{几何断言}：每帧必须恰好 1 次 tuser、720 次 tlast、921,600 个像素
  \item \textbf{帧转储比对}：写出的 \texttt{tpg\_act\_0.hex} 与 Python 生成的
        \texttt{tpg\_exp.hex} 逐像素一致
\end{enumerate}

Python 侧的期望帧用\textbf{同一个颜色表}生成（\texttt{gen\_vectors.py::tpg\_frame}），
保证"设计意图"与"RTL 实现"完全对齐——这是自动化验证的前提。

\subsection{tb\_cdc4phase 的时序陷阱}

\begin{Verbatim}[frame=single,fontsize=\small]
// 正确的驱动时序（修正同沿竞态后）
data_v = WIDTH'($urandom());
s_data <= data_v;
exp_q[wp] <= data_v;
s_load <= 1'b1;
@(posedge clk_s);
s_load <= 1'b0;
@(posedge clk_s);            // 等一拍，让 busy 有条件拉高
wait (s_busy === 1'b1);      // 显式等 busy 拉高（受理）
wait (s_busy === 1'b0);      // 显式等握手完成
sent = sent + 1;
\end{Verbatim}

% ============================================================
\section{验证结果}

\subsection{回归输出}

\begin{Verbatim}[frame=single,fontsize=\small]
$ python tools/run_all.py

--- tb_ser96to24 ...
    ==> PASS
--- tb_tpg ...
    ==> PASS
--- tb_pulse_sync ...
    ==> PASS
--- tb_cdc4phase ...
    ==> PASS
...
==== 汇总: 18/18 PASS ====   ← 当前仓库全量（Step1 阶段为 15 项）
\end{Verbatim}

\subsection{Step1 四项测试细节}

\begin{center}
\renewcommand{\arraystretch}{1.25}
\begin{tabular}{lll}
\toprule
\textbf{测试} & \textbf{验证内容} & \textbf{结果} \\
\midrule
tb\_ser96to24 & 随机空拍 + 3 帧，队列自检（数据/sof/eol 逐像素） & PASS \\
tb\_tpg & 几何计数断言 + 921,600 像素逐像素比对 & PASS \\
tb\_pulse\_sync & 异步双钟 100 脉冲 1:1 传递 & PASS \\
tb\_cdc4phase & 异步双钟 200 个随机 23bit 数据完整性 & PASS \\
\bottomrule
\end{tabular}
\end{center}

\subsection{Step1 阶段抓到的 5 个真缺陷}

\begin{center}
\renewcommand{\arraystretch}{1.25}
\begin{tabular}{clp{7.5cm}}
\toprule
\textbf{\#} & \textbf{现象} & \textbf{根因与修复} \\
\midrule
1 & 串化器溢出（cnt=13） & 源满速发 beat（应为 1/4 占空）$\to$ 全体加 4 分频 \\
2 & CDC 隔包丢 & TB \texttt{wait(!busy)} 同沿竞态 $\to$ 显式等握手完成 \\
3 & TPG 无输出 & \texttt{pres} 漏复位（上电 X）$\to$ 补复位分支 \\
4 & flags 全毁 & hex 字段 2 字符超过 2bit 向量限制 $\to$ 改单字符 \\
5 & TB 编译错 & 缺 \texttt{`timescale} 指令 $\to$ 补上 \\
\bottomrule
\end{tabular}
\end{center}

\begin{okbox}{这 5 个缺陷的价值}
全部 5 个都在\textbf{仿真阶段}发现并修复，没有一个留到上板。\\[4pt]
这正是"金标准 + 逐位比对 + 严格 TB"方法论的价值——
上板调试一次的成本（综合 20 分钟 + 下载 + 示波器测量 + 重新综合）
远高于在仿真里跑几十秒定位问题。
\end{okbox}

% ============================================================
\section{上板挂载说明}

\subsection{在官方例程中的位置}

Step1 的模块\textbf{不改变}官方 ISP$\to$DDR$\to$HDMI 主线，是纯旁路分支：

\begin{Verbatim}[frame=single,fontsize=\small]
isp_top 输出（96bit AXI-S）
    │
    ├──> data_96bit_to_128bit ──> DDR2 ──> HDMI（主线，原封不动）
    │
    └──> detect_top（视觉分支）
             │
             ├── ser96to24 ──> 1ppc ──> HSV/阈值/平滑/质心（Step2）
             │
             └── 直通输出（与输入逐拍一致，TB 已断言 passthru_mism=0）
\end{Verbatim}

\subsection{挂载步骤（给上板组）}

\begin{enumerate}
  \item 在 \texttt{design\_top\_wrapper.v} 中找到 \texttt{isp\_top} 的实例
  \item 把它的输出（\texttt{O\_tdata / O\_tvalid / O\_tuser / O\_tlast}）
        \textbf{并线}接到 \texttt{detect\_top} 的输入（并线，不是串联）
  \item \texttt{detect\_top} 的 \texttt{O\_tready} 悬空或接常量 1（模块内部已固定为 1）
  \item 综合：确认无新警告
  \item 上板：用 ChipWatcher 抓 \texttt{t\_valid} 与 \texttt{t\_tick}，确认每帧有脉冲
\end{enumerate}

\subsection{资源占用}

\begin{center}
\renewcommand{\arraystretch}{1.25}
\begin{tabular}{lll}
\toprule
\textbf{模块} & \textbf{资源估算} & \textbf{说明} \\
\midrule
ser96to24 & $<$ 200 LUT + 8$\times$26bit 分布式 RAM & FIFO 很小 \\
tpg\_720p & $<$ 300 LUT & 只有计数器与比较器 \\
pulse\_sync & $<$ 10 LUT & 4 个触发器 \\
cdc\_4phase & $<$ 100 LUT & 握手逻辑 \\
\midrule
\textbf{Step1 合计} & \textbf{$<$ 610 LUT} & 约占 PH1P35 的 1.6\% \\
\bottomrule
\end{tabular}
\end{center}

\begin{tipbox}{为什么资源这么少}
因为 Step1 只做"搬运"——没有乘法、没有除法、没有大缓存。\\[4pt]
真正的资源消耗在后面：
HSV 查表（256$\times$12 ROM）、行缓存（2$\times$1280bit）、
运动检测帧缓存（115KB ERAM）。
\end{tipbox}


% ============================================================
\section{总结与下一步}

\subsection{Step1 学到了什么}

\begin{itemize}
  \item Verilog 最小语法集（8 个结构 + 3 条铁律）
  \item 视频流三路标（valid/sof/eol）与帧结构
  \item FIFO 的原理和实现（写指针/读指针/计数器）
  \item 串化器的思路（4 像素一包 → 1 像素一拍）
  \item 测试图案发生器的思路（计数器驱动查色表）
  \item CDC 的必要性和两种安全桥
\end{itemize}

\subsection{下一步学什么}

\begin{itemize}
  \item \textbf{Step2 颜色识别与质心}：在像素流上做实时图像处理
  \begin{itemize}
    \item RGB→HSV 颜色空间变换（查表+移位代替除法）
    \item 阈值分割（"这个像素是不是目标颜色？"）
    \item 3×3 平滑（去噪）
    \item 质心计算（"目标在画面哪里？"）
  \end{itemize}
  \item 对应文档：\texttt{03\_vision/step2/02\_技术文档/}
  \item 对应金标准：\texttt{tools/golden\_model.py}
\end{itemize}

\end{document}
